
\assignmenttitle
    {LAB ASSIGNMENT - 02}
    {UNIX SYSCALLS}

\aim{To implement C++ programs using different UNIX system calls and demonstrate their purpose.}

\appr{Computer System with LINUX/UNIX OS and a C++ compiler installed.}

\sectionlabel{Theory:}
A system call is a function provided by the OS through which a user program requests a service from the kernel. It acts as a bridge between user applications and the OS, allowing programs to perform tasks such as creating files, reading data, managing processes, and accessing hardware resources without directly accessing sensitive resources.

\sectionlabel{Process System Calls:}
\begin{itemize}
\item \textbf{fork:} Creates a new child process by duplicating the calling parent process.
\item \textbf{getpid:} Returns the process ID (PID) of the calling process.
\item \textbf{wait:} Suspends execution of the calling process until one of its child processes terminates.
\item \textbf{exec:} Replaces the current process with a new process image from a specified executable file.
\item \textbf{exit:} Terminates the calling process normally and returns a status code to the parent.
\end{itemize}

\sectionlabel{Directory System Calls:}
\begin{itemize}
\item \textbf{opendir:} Opens a directory stream for reading directory entries.
\item \textbf{readdir:} Reads the next directory entry from an open directory stream.
\item \textbf{stat:} Retrieves information about a file such as size, permissions, and modification time.
\item \textbf{close:} Closes a file and releases the associated system resources.
\end{itemize}

\programs

\program{
    2.1. Program 1 - Process System Calls (fork, getpid, wait, exec, exit):
}

\codelabel

\begin{code}
#include <unistd.h>
#include <iostream>
#include <sys/wait.h>

int main() {
    pid_t pid = fork();
    if (pid == 0) {
        std::cout << "Child PID: " << getpid() << std::endl;
        execlp("ls", "ls", NULL);
        exit(1);
    } else {
        wait(NULL);
        std::cout << "Parent PID: " << getpid() << std::endl;
    }
}
\end{code}

\outputlabel
\outputimage[0.50\linewidth]{images/OS/syscalls1.png}

\program{
    2.2 Program 2 - Directory System Calls (opendir, readdir, stat, close):
}

\codelabel

\begin{code}
#include <iostream>
#include <fcntl.h>
#include <unistd.h>
#include <sys/stat.h>
#include <dirent.h>

int main() {
    struct stat fileInfo;
    stat(".", &fileInfo);
    int fd = open(".", O_RDONLY);
    close(fd);

    DIR *dir = opendir(".");
    struct dirent *entry = readdir(dir);

    while (entry != NULL) {
        std::cout << entry->d_name << std::endl;
        entry = readdir(dir);
    }
    closedir(dir);
}
\end{code}

\outputlabel
\outputimage[0.50\linewidth]{images/OS/syscalls2.png}

\concl{The programs demonstrated the use of different system calls to carry out different functions successfully.}
